\documentclass{article}
\usepackage{mathtools} 
\usepackage{fontspec}
\usepackage[UTF8]{ctex}
\usepackage{amsthm}
\usepackage{mdframed}
\usepackage{xcolor}
\usepackage{amssymb}
\usepackage{amsmath}


% 定义新的带灰色背景的说明环境 zremark
\newmdtheoremenv[
  backgroundcolor=gray!10,
  % 边框与背景一致，边框线会消失
  linecolor=gray!10
]{zremark}{说明}

% 通用矩阵命令: \flexmatrix{矩阵名}{元素符号}{行数}{列数}
\newcommand{\flexmatrix}[4]{
  \[
  #1 = \begin{pmatrix}
    #2_{11}     & #2_{12}     & \cdots & #2_{1#4}   \\
    #2_{21}     & #2_{22}     & \cdots & #2_{2#4}   \\
    \vdots      & \vdots      & \ddots & \vdots     \\
    #2_{#31}    & #2_{#32}    & \cdots & #2_{#3#4}
  \end{pmatrix}
  \]
}

% 简化版命令（默认矩阵名为A，元素符号为a）: \quickmatrix{行数}{列数}
\newcommand{\quickmatrix}[2]{\flexmatrix{A}{a}{#1}{#2}}

\begin{document}
\title{第三章总结}
\author{张志聪}
\maketitle

\begin{zremark}
  相比数列求极限，函数求极限一些额外的方法。
\end{zremark}

\begin{itemize}
  \item 利用左右极限相等求极限，这种方式主要应用于分段函数。
  \item 利用数列极限和函数的单调性。如： \S2 例4
\end{itemize}

\begin{zremark}
  利用已经存在的极限，计算遇到的新函数的极限。
  变量替换的原理，运用了一下命题：

  \begin{align*}
    \lim\limits_{x \to x_0} f(x) = A
  \end{align*}
  如果
  \begin{align*}
    \lim\limits_{t \to t_0} \phi(t) = x_0
  \end{align*}
  于是我们有
  \begin{align*}
    \lim\limits_{t \to t_0} (f \circ \phi)(t) = A
  \end{align*}
\end{zremark}

\textbf{证明：}
可以直接使用极限$\epsilon - \delta$定义证明。

\begin{zremark}
  $\circledast$
  常用等价无穷小：

  与$x$等价：
  \begin{align*}
    x \to 0 :       \\
    sin x \sim x    \\
    tan x \sim x    \\
    arctan x \sim x \\
    arcsin x \sim x \\
    e^x - 1 \sim x  \\
    ln(1 + x) \sim x
  \end{align*}
  其他
  \begin{align*}
    x \to 0 :                    \\
    1 - cosx \sim \frac{1}{2}x^2 \\
    (1 + x)^a - 1 \sim a x \ (a \neq 0)
  \end{align*}
\end{zremark}

\end{document}